%% en/sections/06_esiti.tex — Outcomes (full EN)

\chapter{The Outcomes}
\markboth{The Outcomes}{The Outcomes}

\lettrine[lines=3]{T}{he} evening of 30 March 1282 is inevitable.
The question is not \textit{whether} something will happen --- but
\textit{what}, and in what form. The scoring system measures how
effectively the factions have prepared, organised, or obstructed
the night that is coming.

\section{The Scoring System}

\subsection{Calculating the Total Score}

At the end of the session, the coordinating GM collects all four
faction scores:

\begin{center}
\begin{tabular}{lc}
\toprule
\cinzel Faction & \cinzel Score (0--4) \\
\midrule
Conspirators  & \hspace{2cm} \\[4pt]
Barons        & \hspace{2cm} \\[4pt]
Clergy        & \hspace{2cm} \\[4pt]
People        & \hspace{2cm} \\
\midrule
\textbf{Total} & \textbf{\hspace{2cm} / 16} \\
\bottomrule
\end{tabular}
\end{center}

\nota[For the Coordinating GM]{
  The score is calculated after the session ends, before revealing the
  outcome to players. The reveal is a narrative moment: do not simply
  read the text. Adapt it to what actually happened at your tables.
}

\secsep

\section{The Outcomes}

\subsection{0--4 points --- \textit{The Night of Silence}}
\textcolor{oro}{\cinzel\large Conspiracy Broken}

\filetto

The conspiracy has been broken, or was never organised enough to
function. Jean de Saint-Rémy moved first: pre-emptive arrests, curfew,
port sealed. The Santo Spirito procession takes place under heavy
military surveillance.

Nothing happens on 30 March.

But the seeds have been planted. The injustices remain. The anger
remains. The Sicilian Vespers are postponed --- not cancelled.
The history of the Mediterranean will be different. How different,
no one can know.

\textbf{Faction consequences:}
\begin{itemize}[nosep]
  \item \textit{Conspirators} --- the network is partially or totally
    dismantled. Giovanni da Procida must start again from scratch.
  \item \textit{Barons} --- those who had taken a position are compromised.
    Those who had kept a foot in both camps congratulate themselves.
  \item \textit{Clergy} --- the Church remains in place, neither more
    nor less respected than before.
  \item \textit{People} --- repression falls on the quarters. Ordinary
    people pay the price of what others attempted.
\end{itemize}

\nota[Suggested Epilogue]{
  Read players a short text set one year later: the uprising broke out
  elsewhere, in a different way, with unforeseeable consequences.
  Their characters are alive --- or perhaps not. But history moved
  on without them.
}

\secsep

\subsection{5--8 points --- \textit{The Spark Extinguished}}
\textcolor{oro}{\cinzel\large Partial Uprising, Suppressed}

\filetto

The Santo Spirito procession takes place. There is an incident ---
a search that went too far, a soldier who does not know when to stop.
A brawl breaks out. In some quarters the brawl becomes violence. In
others, the authorities manage to contain it.

The uprising exists --- but it is disorganised, localised, and suppressed
within a few days. Dozens of civilian dead. Charles of Anjou uses the
episode to justify harsher repression. Peter of Aragon does not intervene:
he was not ready, and the window has closed.

\textbf{Faction consequences:}
\begin{itemize}[nosep]
  \item \textit{Conspirators} --- the plan was there, but not robust
    enough. Some fled. Some were captured.
  \item \textit{Barons} --- those who bet on the uprising lose everything.
    Those who kept a foot in both camps survive, at a price.
  \item \textit{Clergy} --- an ecclesiastical inquiry is opened.
    The Archbishop is summoned to answer to Rome.
  \item \textit{People} --- the quarters mourn their dead. The memory
    of this night will feed the next uprising.
\end{itemize}

\secsep

\subsection{9--12 points --- \textit{The Vespers}}
\textcolor{oro}{\cinzel\large Full Uprising}

\filetto

The Santo Spirito procession becomes the detonator. The incident is
real, the moment is right, and the network built by the factions holds.
Within hours all Palermo is in revolt.

French soldiers are driven out or killed. Jean de Saint-Rémy manages
to flee from the port before it is sealed. By the morning of 31 March,
Palermo is free for the first time in sixteen years.

This is what happened historically --- but in this version, the players
built it. Some things went differently. Some people are alive who
historically were not, or the reverse.

\textbf{Faction consequences:}
\begin{itemize}[nosep]
  \item \textit{Conspirators} --- the network held. Da Procida and
    Peter of Aragon have their opportunity. What they do with it is
    another story.
  \item \textit{Barons} --- those who bet correctly have won --- for now.
    The war has just begun.
  \item \textit{Clergy} --- the bells of Santo Spirito rang at the
    right moment. The Sicilian Church has new autonomy --- and new
    responsibilities.
  \item \textit{People} --- the victory belongs to those who were in
    the square. But the decisions about what happens next will not be
    theirs.
\end{itemize}

\nota[Narrative Variants]{
  At 9--10 points: the uprising succeeds, but with more violence and
  less control than planned. At 11--12: the uprising succeeds close
  to the plan. Adapt the narration accordingly.
}

\secsep

\subsection{13--16 points --- \textit{The Insurrection}}
\textcolor{oro}{\cinzel\large Total, Coordinated Uprising}

\filetto

The night of 30 March is not an uprising --- it is an insurrection.
The four factions worked with enough coordination to transform an
incident into a planned revolution. Not just Palermo: the channels
opened to Messina, Corleone, Trapani carried the uprising across
Sicily with a speed no one expected --- not even the conspirators.

Peter of Aragon lands earlier than planned. Charles's expedition against
Byzantium is permanently cancelled. The Mediterranean is changed.

\textbf{Faction consequences:}
\begin{itemize}[nosep]
  \item \textit{Conspirators} --- Giovanni da Procida is triumphant.
    But triumph has a price: who decides what happens now is Peter of
    Aragon, not the Sicilians.
  \item \textit{Barons} --- the houses that chose correctly are the
    new pillars of the kingdom. Heavy responsibility.
  \item \textit{Clergy} --- the Sicilian Church has gained unprecedented
    autonomy --- and a debt to Aragon.
  \item \textit{People} --- they won. Their names will not be in the
    chronicles. But they know it.
\end{itemize}

\nota[Design Note --- The Outcome Is Not a Reward]{
  None of the outcomes is intrinsically ``good'' or ``bad'' in a moral
  sense. Even total insurrection has a cost --- the violence of that
  night is real, the consequences are unpredictable, and Peter of Aragon
  is not a disinterested liberator. The game does not reward those who
  made the ``right'' choices: it records what happened, and lets players
  carry the weight.
}

\secsep

\section{Special Faction Outcomes}

Independently of the total score, certain special conditions modify or
add elements to the final narration:

\begin{itemize}
  \item \textbf{Conspirators at 0} --- the network is completely dismantled,
    even if the uprising succeeds: it broke out without them, against their
    will or in their absence.
  \item \textbf{Barons at 4} --- the Sicilian baronage enters the new era
    united: rare and powerful. A Baron character gains a title in the new
    Sicily.
  \item \textbf{Clergy at 4 and Sacristan alive} --- the bells of Santo
    Spirito ring as dawn breaks. It is the sound of victory. It is the
    sound they will remember.
  \item \textbf{People at 4} --- the popular quarters suffer no retaliation
    that night. All People characters and their loved ones survive.
  \item \textbf{De Saint-Rémy Clock full} --- even in a successful uprising,
    the repression claimed lives. Read the names.
\end{itemize}

\filettodoppio
